$k$ is $1$-indexed, assuming dummy on both end. \texttt{insert} inserts after $k$-th vertex (front if $k = 0$).

\complexity{$\mathcal O(Q \log N)$, amortized.}
\begin{minted}{cpp}
const int INF = 1e9;
struct Node{
	int p = -1, l = -1, r = -1; bool flp = 0; int val = 0;
	int cnt = 1; ll sum = 0; int mn = INF, mx = -INF;
}; vector<Node> tree; int root = -1;
Node init(int x){
	Node p; p.cnt = 1; p.val = x;
	p.sum = p.mn = p.mx = x; return p;
} void put(int ni, const Node p){
	tree[ni].cnt = p.cnt; tree[ni].sum = p.sum; tree[ni].mn = p.mn; tree[ni].mx = p.mx;
}
Node mrg(int li, int ri){
	const Node &l = tree[li], &r = tree[ri];
	Node p; p.cnt = l.cnt + r.cnt;
	p.sum = l.sum + r.sum; p.mn = min(l.mn, r.mn); p.mx = max(l.mx, r.mx);
	return p;
}
void propagate(int ni){
	if (ni == -1){ return; } if (tree[ni].flp == 0){ return; }
	swap(tree[ni].l, tree[ni].r);
	if (tree[ni].l != -1){ tree[tree[ni].l].flp ^= 1; }
	if (tree[ni].r != -1){ tree[tree[ni].r].flp ^= 1; }
	tree[ni].flp = 0;
}
void update(int ni){
	put(ni, init(tree[ni].val));
	if (tree[ni].l != -1){ propagate(tree[ni].l); put(ni, mrg(tree[ni].l, ni)); }
	if (tree[ni].r != -1){ propagate(tree[ni].r); put(ni, mrg(ni, tree[ni].r)); }
} inline void update(int ni, ll val){ tree[ni].val = val; return update(ni); }
void rotate(int ni){
	propagate(tree[ni].p); propagate(ni);
	int p1 = tree[ni].p; if (p1 == -1){ return; }
	int p2 = tree[p1].p; if (p2 != -1){
		(tree[p2].l == p1 ? tree[p2].l : tree[p2].r) = ni;
		tree[ni].p = p2;
	} else{ root = ni; tree[ni].p = -1; }
	int p3 = -1; if (tree[p1].l == ni){
		tree[p1].l = p3 = tree[ni].r; tree[ni].r = p1;
	} else{
		tree[p1].r = p3 = tree[ni].l; tree[ni].l = p1;
	} tree[p1].p = ni; if (p3 != -1){ tree[p3].p = p1; }
	update(p1); update(ni);
}
int splay(int ni, int top = -1){
	while (tree[ni].p != top){
		propagate(tree[tree[ni].p].p); propagate(tree[ni].p); propagate(ni);
		int p1 = tree[ni].p; int p2 = tree[p1].p;
		if (p2 != top){ rotate((tree[p2].l==p1) == (tree[p1].l==ni) ? p1 : ni); }
		rotate(ni);
	} propagate(ni); return ni;
}
int kth(int k){
	int ni = root; while (1){
		propagate(ni);
		int cnt = (tree[ni].l == -1 ? 0 : tree[tree[ni].l].cnt);
		if (cnt == k){ break; }
		if (k < cnt){ ni = tree[ni].l; }
		else{ k -= cnt+1; ni = tree[ni].r; }
	} return splay(ni);
}
int segment(int st, int ed){
	kth(ed+1); int r = root; kth(st-1); splay(r, root); return tree[tree[root].r].l;
}
void insert(int k, int val){
	if (root == -1){ root = tree.size(); tree.push_back(init(val)); return; }
	int ni = root; while (1){
		int cnt = (tree[ni].l == -1 ? 0 : tree[tree[ni].l].cnt);
		if (k < cnt){
			if (tree[ni].l == -1){
				tree[ni].l = tree.size(); tree.push_back(init(val));
				tree[tree[ni].l].p = ni; ni = tree[ni].l; break;
			} else{ ni = tree[ni].l; }
		} else{
			if (tree[ni].r == -1){
				tree[ni].r = tree.size(); tree.push_back(init(val));
				tree[tree[ni].r].p = ni; ni = tree[ni].r; break;
			} else{ k -= cnt+1; ni = tree[ni].r; }
		}
	} splay(ni);
}
void erase(int k){
	segment(k, k); tree[tree[root].r].l = -1; splay(tree[root].r);
}

void rev(int st, int ed){
	if (st > ed){ return; }
	int ni = segment(st, ed); tree[ni].flp ^= 1; propagate(ni);
}
\end{minted}